\section{Teoría Greedy}

''La codicia...es buena. La codicia es correcta. La codicia funciona.'' - Michael Douglas (película Wall Street 1987).
Codicia = Greed. 

Es difícil de explicar bien qué es esta terminología. Pero estamos hablando de un razonamiento \textit{sin} visión general, sin la visión del ''Big Picture'' del problema o planteamiento.

Se asocia a la \textit{avaricia} de la razón por la acción tomada. Y ahora piensen en Sophia, que quiere ganar siempre, que siempre agarra las monedas que a ella le convenga con tal de ganar en todo momento. Eso es avaricia.

Esa acción, ese pensamiento, esa \textit{secuencia} de acciones se puede denominar como una acción Greedy. Siempre, en el momento en el que estoy y donde estoy, busco lo mejor para mí. Sin importar lo que viene después.

Ahora más formalizado, se aplica una regla sencilla que nos permite obtener el/un \textbf{óptimo local} a mi \textbf{estado actual}. Aplicamos iterativamente esa regla, esperando que nos lleve al óptimo global (que sería, la mejor solución).

Un algoritmo es greedy cuando tenemos dicha regla sencilla que sólo mira el estado actual y que busca un óptimo local. Si no podemos enunciar cuál es el óptimo local, el algoritmo no es greedy o lo \textit{estamos pensando mal}.

Pensar de esta manera, como veníamos diciendo sin la visión general del problema, \textit{esperando} que nos lleve al óptimo global, los algoritmos Greedy nos pueden indicar varias cosas:
\begin{itemize}
    \item No siempre dan el resultado óptimo.
    \item Demostrar que dan el resultado óptimo puede no ser fácil.
    \item Son intuitivos de pensar. Son fáciles de entender.
    \item Suelen funcionar rápido.
    \item Para problemas complejos pueden ser buenas aproximaciones.
\end{itemize}

Recordar la regla de: Sucesión de óptimos locales -> óptimo global. Esa misma nos va a permitir explicar el por qué un algoritmo es greedy. 


\subsection{Regla Greedy buscada}
. explico que es un algoritmo greedy

. explico la regla greedy aplicada a este ejercicio

. explico el estado actual, el optimo local, el optimo global y la sucesion

. digo que la demostracion va a estar al final
\subsection{Pensamiento Greedy}
JEE 

Filosófico:
algoritmo greedy = Liberalismo?

Consecuencialismo / Utilitarismo vs Deontologismo.


Frase Greedy Kleinberg Tardos.